\documentclass[10pt]{article}
\usepackage[T1]{fontenc}
\usepackage[utf8]{inputenc}
\usepackage[french]{babel}
\usepackage{amsmath,amssymb,mathrsfs,mathpazo}
\usepackage{geometry}
\geometry{a4paper,margin=2cm}
\usepackage{graphicx,array,tabularx,tabularray,colortbl,multirow,multicol,booktabs,enumitem,pifont,ulem,fancybox,diagbox,cancel}
\usepackage{pstricks,pstricks-add,pst-plot,pst-tree,pst-node,pst-text,pst-func,pst-math,pst-bspline,pst-3dplot,pst-eucl,pst-solides3d}
\let\pscalculate\undefined
\usepackage{pst-bezier}
\usepackage{tikz,pgfplots,tkz-tab}
\usetikzlibrary{arrows,arrows.meta,patterns,calc,positioning,trees,decorations.markings,shapes.geometric,shapes.misc,shapes.arrows,intersections,decorations.pathmorphing,decorations.pathreplacing,backgrounds,shadows,babel}
\pgfplotsset{compat=1.18}
\usepgfplotslibrary{fillbetween}
\usepackage[np]{numprint}
\usepackage{listings,makecell,pas-tableur,twoopt}
\let\barre\undefined
\usepackage[pstricks]{tablvar}
\definecolor{darkgreen}{RGB}{0,100,0}
\providecommand{\red}{\color{red}}
\providecommand{\blue}{\color{blue}}
\providecommand{\green}{\color{green}}
\providecommand{\black}{\color{black}}
\providecommand{\white}{\color{white}}
\providecommand{\gray}{\color{gray}}
\providecommand{\cyan}{\color{cyan}}
\providecommand{\magenta}{\color{magenta}}
\providecommand{\yellow}{\color{yellow}}
\providecommand{\violet}{\color{violet}}
\providecommand{\orange}{\color{orange}}
\providecommand{\euro}{\texteuro}
\providecommand{\EUR}{\texteuro}
\providecommand{\ieme}{\textsuperscript{e}}
\providecommand{\degres}{\ensuremath{^\circ}}
\providecommand{\degree}{\ensuremath{^\circ}}
\providecommand{\vv}[1]{\overrightarrow{#1}}
\providecommand{\Coord}[1]{#1}
\providecommand{\decosix}{}
\providecommand{\starredbullet}{\ensuremath{\bullet}}
\pagestyle{empty}
\setlength{\parindent}{0pt}
\usepackage{adjustbox}\DeclareUnicodeCharacter{2236}{:}\DeclareUnicodeCharacter{2190}{\ensuremath{\leftarrow}}
\providecommand{\R}{\mathbb{R}}
\providecommand{\N}{\mathbb{N}}
\providecommand{\Z}{\mathbb{Z}}
\providecommand{\Q}{\mathbb{Q}}
\providecommand{\D}{\mathbb{D}}
\providecommand{\C}{\mathbb{C}}
\providecommand{\e}{{\mathrm{e}}}
\providecommand{\vect}[1]{\overrightarrow{#1}}
\providecommand{\vectt}[1]{\overrightarrow{\mathrm{#1}}}
\providecommand{\vvt}[1]{\overrightarrow{\mathrm{#1}}}
\providecommand{\np}[2][]{{#2~#1}}
\providecommand{\vv}[1]{{\overrightarrow{#1}}}
\renewcommand{\d}{{\mathrm{d}}}
\providecommand{\degres}{{^\circ}}
\providecommand{\degre}{{^\circ}}
\providecommand{\blue}{\color{blue}}
\providecommand{\red}{\color{red}}
\providecommand{\enskip}{\quad}
\providecommand{\llbracket}{\lbrack\!\lbrack}
\providecommand{\rrbracket}{\rbrack\!\rbrack}
\providecommand{\euro}{{\text{€}}}
\providecommand{\no}{{\text{n°}}}
\providecommand{\interval}[2]{{[#1;#2]}}
\providecommand{\intervFO}[2]{{[#1;#2[}}
\providecommand{\dotfill}{\ldots}

\providecommand{\ds}{\displaystyle}
\providecommand{\textding}[1]{#1}
\providecommand{\Oij}{\(\mathrm{O};\vec{\imath},\vec{\jmath}\)}
\providecommand{\Oijk}{\(\mathrm{O};\vec{\imath},\vec{\jmath},\vec{k}\)}
\providecommand{\Ouv}{\(\mathrm{O};\vec{u},\vec{v}\)}
\providecommand{\intervFF}[2]{\ensuremath{[#1;#2]}}
\renewcommand{\euro}{\ifmmode\text{\texteuro}\else\texteuro\fi}
\DeclareUnicodeCharacter{00B0}{\ifmmode^{\circ}\else\textdegree\fi}
\DeclareUnicodeCharacter{20AC}{\ifmmode\text{\texteuro}\else\texteuro\fi}
% Énoncés : Éducation nationale. Transcription et corrigés : APMEP, Denis Vergès et contributeurs. Découpage et classement : CoopMaths, Sébastien Lozano, Éric Elter et contributeurs. Adaptation de la présentation : Maths Explorer.
% CC BY-SA : https://coopmaths.fr/www/about/
% https://coopmaths.fr/alea/static/bac/2024/tex/bac_2024_01_sujet0_5_cor.tex
\begin{document}
\noindent\textit{APMEP / CoopMaths. CC BY-SA. Présentation adaptée par Maths Explorer.}\par\medskip
\begin{enumerate}
	\item Sur l'intervalle $\left [0~;~2 \pi\strut \right ]$, l'équation  $\sin (x) = 0,1$ admet :

	\begin{tabularx}{\linewidth}{*{2}{>{\bfseries\arraybackslash}p{5mm}@{}X}}
		a.& zéro solution&
		b.& une solution\\
		c.& deux solutions&
		d.& quatre solutions\\
	\end{tabularx}

On peut par exemple dire que la courbe représentant la fonction sinus et la droite d'équation $y=0,1$ ont deux points d'intersection sur l'intervalle $\left [ 0\;;\; 2\pi\strut \right ]$.

\begin{center}
\psset{xunit=1cm,yunit=1cm}
\begin{pspicture}(-1,-1.5)(7,1.5)
\psaxes[trigLabels=true,trigLabelBase=2,dx=\psPiH,xunit=\psPi]{->}(0,0)(0,-1.5)(2.2,1.5)
\psplot[linewidth=1.2pt,linecolor=red,plotpoints=5000]{0}{6.2832}{x RadtoDeg sin}
\psline[linecolor=blue](0,0.1)(6.2832,0.1)
\end{pspicture}
\end{center}

\hfill\textbf{Réponse c.}

	\item On considère la fonction $f$ définie sur l'intervalle $\left [0~;~\pi\strut\right ]$ par  $f(x) = x + \sin (x).$


	\begin{tabularx}{\linewidth}{{>{\bfseries\arraybackslash}p{5mm}@{}X}}
		a.& La fonction $f$ est convexe sur l'intervalle $[0~;~\pi]$\\
		b.& La fonction $f$ est concave sur l'intervalle $[0~;~\pi]$\\
		c.& La fonction $f$ admet sur l'intervalle $[0~;~\pi]$ un unique point d'inflexion\\
		d.& La fonction $f$ admet sur l'intervalle $[0~;~\pi]$ exactement deux points d'inflexion\\
	\end{tabularx}

$f'(x)=1+\cos(x)$ et $f''(x)=-\sin(x)\leqslant 0$ sur $\left [0~;~\pi\strut\right ]$, donc la fonction $f$ est concave sur cet intervalle.

\hfill\textbf{Réponse b.}

	\item Une urne contient cinquante boules numérotées de 1 à 50. On tire successivement trois boules dans cette urne, \textbf{sans remise}. On appelle \og{} tirage \fg{} la liste non ordonnée des numéros des trois boules tirées.
Quel est le nombre de tirages possibles, \textbf{sans tenir compte de l'ordre des numéros}?

	\begin{tabularx}{\linewidth}{*{4}{>{\bfseries\arraybackslash}p{5mm}@{}X}}
		a.& $50^{3}$&
		b.& $1 \times 2 \times 3$&
		c.& $50 \times 49 \times 48$&
		d.& $\dfrac{50 \times 49 \times 48}{1 \times 2 \times 3}$\\
	\end{tabularx}

$\ds\binom{50}{3}=\dfrac{50\times 49 \times 48}{1 \times 2 \times 3}$
\hfill\textbf{Réponse d.}

	\item On effectue dix lancers d'une pièce de monnaie. Le résultat d'un lancer est \og{}pile\fg{} ou \og{}face\fg{}. On note la liste ordonnée des dix résultats.\\
Quel est le nombre de listes ordonnées possibles?

	\begin{tabularx}{\linewidth}{*{2}{>{\bfseries\arraybackslash}p{5mm}@{}X}}
		a.& $2 \times 10$&
		b.& $2^{10}$\\
		c.& $1 \times 2 \times 3 \times \cdots \times 10$&
		d.& $\dfrac{1 \times 2 \times 3 \times \cdots \times 10}{1 \times 2}$\\
	\end{tabularx}

Il y a 2 résultats possibles si on effectue 1 lancer, $2^2$ résultats possibles si on effectue 2 lancers, etc., $2^{10}$ résultats possibles si on effectue 10 lancers.

\hfill\textbf{Réponse b.}

	\item On effectue $n$ lancers d'une pièce de monnaie équilibrée. Le résultat d'un lancer est \og{}pile\fg{} ou \og{}face\fg{}. On considère la liste ordonnée des $n$ résultats.\\
Quelle est la probabilité d'obtenir au plus deux fois \og{} pile \fg{} dans cette liste ?

	\begin{tabularx}{\linewidth}{*{2}{>{\bfseries\arraybackslash}p{5mm}@{}X}}
		a.& $\dfrac{n(n-1)}{2}$&
		b.& $\dfrac{n(n-1)}{2} \times\left(\dfrac{1}{2}\right)^{n}$\\
		c.& $1 + n + \dfrac{n(n-1)}{2}$&
		d.& $\left(1 + n + \dfrac{n(n-1)}{2}\right) \times\left(\dfrac{1}{2}\right)^{n}$
	\end{tabularx}
	
La variable aléatoire $X$ qui donne le nombre de résultats \og pile \fg{} sur $n$ lancers suit la loi binomiale de paramètres $n$ et $\dfrac{1}{2}$.

On cherche $P(X\leqslant 2)$ soit $P(X=0)+P(X=1)+P(X=2)$.
 
\textbullet~~$P(X=0)=\ds\binom{n}{0}\left (\dfrac{1}{2}\right )^0 \left (1-\dfrac{1}{2}\right )^{n-0} = \left (\dfrac{1}{2}\right )^n$

\textbullet~~$P(X=1)=\ds\binom{n}{1}\left (\dfrac{1}{2}\right )^1 \left (1-\dfrac{1}{2}\right )^{n-1} = n \left (\dfrac{1}{2}\right )^n$

\textbullet~~$P(X=2)=\ds\binom{n}{2}\left (\dfrac{1}{2}\right )^2 \left (1-\dfrac{1}{2}\right )^{n-2} = \dfrac{n(n-1)}{2} \left (\dfrac{1}{2}\right )^n$

Donc $P(X\leqslant 2)= \left (\dfrac{1}{2}\right )^n + n \left (\dfrac{1}{2}\right )^n + \dfrac{n(n-1)}{2}\left (\dfrac{1}{2}\right )^n = \left (1+n+\dfrac{n(n-1)}{2}\right )\left (\dfrac{1}{2}\right )^n$.

\hfill\textbf{Réponse d.}
\end{enumerate}

\end{document}
